% ═══════════════════════════════════════════════════════════════════════
%  Stratégie 2 --- XAUUSD niveaux x10. Chiffres : macros générées.
%  Les pistes du §3 sont des HYPOTHÈSES, formulées comme telles. Aucune
%  n'est une promesse, aucune n'est testée.
% ═══════════════════════════════════════════════════════════════════════
\section{Motifs de non-déploiement}
\label{sec:x10_non_deploiement}

\lettrine[lines=2, findent=2pt, nindent=0pt]{L}{a} recommandation est de
\textbf{ne pas déployer} la stratégie XAUUSD sur les niveaux x10, ni en
capital réel ni sur un compte de démonstration. Cette section énumère les
motifs, dit ce qui manque encore, et distingue soigneusement ce que les
mesures autorisent à penser de ce qu'elles ne permettent que de supposer.

\subsection{Ce que les mesures établissent}

Six constats, tous mesurés, tous cohérents entre eux.

\begin{enumerate}
  \item \textbf{L'espérance nette est négative et large.}
  \metric{\XTenExpectancyR{}~R} par transaction, sur \XTenTradesIS{}
  transactions, pour un profit factor de \XTenProfitFactor. Le seuil gelé
  demandait une espérance d'au moins $+0{,}10$~R et un profit factor d'au
  moins $1{,}15$.

  \item \textbf{Aucun réglage ne tient.} Les vingt-sept Sharpes nets sont
  négatifs, de \XTenPeakSharpe{} à \XTenWorstGridSharpe.
  \XTenPlateauPassing{} configuration passe le critère de plateau, et la règle
  gelée interdit de se rabattre sur le pic.

  \item \textbf{Aucune période ne tient.} \XTenPositiveYears{} année sur
  \XTenYearsCount{} est positive. La perte est diffuse~:
  \XTenProfitableSessions{} séances gagnantes contre
  \XTenLosingSessions{} perdantes, sans concentration sur un épisode
  particulier.

  \item \textbf{Aucune population de transactions ne tient.} Les quatre
  scénarios, les trois séances, les trois terciles de volatilité et les
  quatre tranches de régime sont tous négatifs. La coupe la moins mauvaise
  --- la tranche de régime intermédiaire --- reste à
  \XTenBestDollarBandR{}~R.

  \item \textbf{Le résultat survit au rééchantillonnage.} L'intervalle de
  confiance à 95~\% du Sharpe, \XTenSharpeCILow{} à \XTenSharpeCIHigh, ne
  contient pas zéro. Le rééchantillonnage de l'ordre des transactions montre
  que la séquence réalisée n'a \emph{pas} été malchanceuse.

  \item \textbf{Le coût n'est pas la cause.} À coût nul extrapolé,
  l'espérance resterait autour de \XTenZeroCostExtrapolation{}~R --- une
  extrapolation indicative, non mesurée, qui suffit à écarter l'explication
  par les frais.
\end{enumerate}

\begin{warningbox}
\textbf{Le \emph{Deflated Sharpe Ratio} vaut \XTenDSR.} Compte tenu des
\XTenTrialsCount{} essais consommés, le meilleur d'entre eux aurait un Sharpe
attendu de \XTenExpectedMaxSharpe{} par pur hasard. Le Sharpe observé est
\XTenSharpeIS. La stratégie ne bat pas le bruit~; elle lui est très
inférieure.
\end{warningbox}

\subsection{Ce qui manque encore, et pourquoi cela ne change rien}

Deux critères sur neuf sont en échec par \emph{non-mesure}~: le portage
MetaTrader~5 et la lecture hors échantillon. La table de décision impose de
les compter comme des échecs, jamais comme « non applicables ».

\begin{description}
  \item[Le portage MetaTrader~5.] Sa campagne est en cours
  (section~\ref{sec:x10_reconciliation}). Deux remarques~: il ne couvrira
  qu'un peu moins de la moitié de l'échantillon, faute d'historique plus
  ancien chez le courtier~; et son modèle de remplissage interpolé en fait un
  \emph{majorant}. Un chiffre d'exécution qui contredirait six critères de
  recherche serait un motif d'enquête sur le portage, pas un motif de
  déploiement.
  \item[La lecture hors échantillon.] Elle n'a pas eu lieu, et la
  section~\ref{sec:x10_hors_echantillon} explique pourquoi c'est la bonne
  décision~: le verdict est acquis, le holdout est un budget non
  renouvelable, et le régime récent est hors de la distribution de sélection.
\end{description}

\subsection{Ce qui pourrait rouvrir le dossier}

Cette sous-section est la plus délicate du rapport, et elle appelle un
avertissement préalable.

\begin{warningbox}
\textbf{Une décomposition \emph{ex post} n'est pas une preuve.} Les pistes
ci-dessous viennent toutes de coupes faites \emph{après} avoir vu les
résultats. Découper un échantillon perdant en sous-groupes produit
mécaniquement des sous-groupes moins perdants que d'autres, même quand aucune
structure n'existe. Aucune de ces pistes n'est testée, aucune n'est
recommandée, et aucune ne doit être lue comme une promesse de rendement.

Chacune relève d'une \textbf{nouvelle spécification et d'un nouveau budget
d'essais}, avec sa propre table de décision gelée avant mesure. Les greffer
sur la spécification actuelle serait exactement la faute que tout le
dispositif de cette phase cherchait à empêcher.
\end{warningbox}

\paragraph{Hypothèse 1 --- une grille adaptative en volatilité.}
La dérive du pas de grille est le fait le plus massif du dossier~: dix
dollars valent \XTenDollarsPerAtrFirst{} \code{ATR} en début d'échantillon et
\XTenDollarsPerAtrLast{} à la fin, et les deux tranches extrêmes de régime
sont les pires. On pourrait supposer qu'une grille dont le pas s'exprimerait
en multiples d'\code{ATR} --- et non en dollars --- présenterait au signal un
objet stable d'une année à l'autre.

\emph{Ce qui s'y oppose.} Ce ne serait plus la stratégie commandée~: le
mandat désigne les niveaux psychologiques de dix dollars, c'est-à-dire des
prix ronds, et un niveau « à 3,2~\code{ATR} du précédent » n'a aucune
existence psychologique. L'hypothèse mérite d'être formulée~; elle change
d'objet d'étude.

\paragraph{Hypothèse 2 --- un scénario unique.}
Les \emph{breakouts} perdent moins que les \emph{reversals}, et le
\emph{Breakout Long} est le moins mauvais des quatre à
\XTenBestScenarioR{}~R. On pourrait supposer qu'une stratégie réduite à cette
seule famille, avec un budget de conception qui lui serait entièrement
consacré, se comporterait différemment.

\emph{Ce qui s'y oppose.} Le \emph{Breakout Long} reste négatif, sur un
nombre de transactions suffisant pour que ce ne soit pas du bruit. Et une
partie de son avantage vient probablement du fait qu'il est le seul, avec son
symétrique, à passer un filtre de contexte --- filtre dont les ablations
montrent qu'il améliore sans redresser.

\paragraph{Hypothèse 3 --- une restriction de séance.}
C'est la coupe la plus contrastée de la campagne~: la séance asiatique porte
\XTenAsiaLossShare{} de la perte pour \XTenAsiaTradeShare{} des transactions,
avec \XTenAsiaExpectancyR{}~R contre \XTenNewYorkExpectancyR{}~R en séance
new-yorkaise. On pourrait supposer qu'une stratégie limitée aux heures de
Londres et de New York éviterait la part la plus coûteuse du flux.

\emph{Ce qui s'y oppose.} D'abord, la séance new-yorkaise \emph{reste
négative}~: retirer l'Asie ne ferait pas passer l'espérance du bon côté de
zéro, cela la rapprocherait seulement. Ensuite, le mandat demande
explicitement la séance entière. Enfin, un précédent du dépôt --- une
composante horaire intraday abandonnée --- avait identifié l'absence de
restriction de séance comme facteur aggravant~; le reprendre ici serait
cohérent, mais cela resterait un choix fait après avoir vu la décomposition.

\paragraph{Hypothèse 4 --- l'extension au VWAP comme filtre.}
C'est la seule variable descriptive de la campagne qui sépare réellement
quelque chose~: les \emph{reversals} pris loin du \code{VWAP} de séance
perdent \XTenExtendedR{}~R contre \XTenNotExtendedR{}~R pour les autres.
La spécification l'avait figée comme tracée et non bloquante, par décision du
mandat.

\emph{Ce qui s'y oppose.} L'écart est net mais les deux branches sont
négatives~: en faire un filtre bloquant ne produirait pas une espérance
positive. Et c'est exactement le type de variable qu'un échantillon perdant
fait apparaître par construction.

\subsection{Ce qui a été appris, indépendamment du verdict}

Un dossier négatif produit tout de même des acquis, et ceux-ci survivront à ce
rapport.

\begin{itemize}
  \item \textbf{Une chaîne de décision complète, spécifiée et testée} pour une
  stratégie intraday sur niveaux, portable sur trois moteurs~: horloge de
  séance, causalité multi-échelles, machine à états, exécution pessimiste,
  contrat de trace événementielle.

  \item \textbf{Une mesure durable de la non-stationnarité de l'or.} Le
  tableau du pas de grille en \code{ATR} n'est pas propre à cette stratégie~:
  il s'applique à toute règle fondée sur une distance en dollars, et il
  documente un régime qui a changé d'un facteur supérieur à quatre en sept
  ans.

  \item \textbf{Un contre-exemple sur le retest.} Le mandat le citait comme
  une confirmation~; la mesure le montre contre-prédictif
  (\XTenRetestR{}~R avec, \XTenNoRetestR{}~R sans). C'est une information
  utile au-delà de ce dossier.

  \item \textbf{Une erreur de méthode documentée.} Le garde-fou gain/risque
  avait été jugé inopérant \emph{par raisonnement}, sur une hypothèse
  d'entrée idéalisée~; le comptage l'a réfuté --- il rejette
  \XTenRRejectLastYear{} des candidats en fin d'échantillon. La leçon est
  conservée dans les notes de recherche.

  \item \textbf{Deux défauts d'outillage identifiés}~: le test de supériorité
  prédictive à orientation inversée et le test descendant qui échoue sur une
  grille sans gagnant. Ils étaient préexistants~; ils sont désormais nommés.
\end{itemize}

\begin{insightbox}
\textbf{Le risque réel de ce chantier n'était pas de mesurer un échec~: c'était
de re-sélectionner après l'avoir mesuré.} Une grille figée à vingt-sept
points, un budget de \XTenTrialsCeiling{} essais dont \XTenTrialsReserve{}
restent inutilisés, des seuils gelés et committés avant la première mesure, un
plateau exigé plutôt qu'un pic, et un holdout qui n'a pas été ouvert. Le
dispositif a tenu~: le verdict de ce rapport est le verdict que la table de
décision prévoyait, appliquée à la lettre.
\end{insightbox}
